\ifajar
  \section{Latihan}
\else
  \section{Bab 6: Testing, Reliability, dan Code Quality}
\fi

Seluruh file yang dipakai pada latihan ini tersedia di
\berkas{sources/chapter06/}, dihitung relatif terhadap akar repositori.
Latihan membutuhkan Docker dengan Compose versi 2, Python 3, Node.js, dan
Google Chrome yang sudah terpasang. Listing penyiapan masuk ke direktori itu
lebih dahulu, dan seluruh perintah sesudahnya dijalankan dari sana.

Seluruh kegiatan pada bab ini hanya menyentuh basis data dan aplikasi milik
sendiri di komputer sendiri, sehingga tidak ada sistem pihak lain yang
dibebani. Hasil tiap mahasiswa akan berbeda, karena waktunya bergantung pada
perangkat keras dan beban lain yang sedang berjalan. Yang dilaporkan adalah
pola perbandingannya, bukan angka persisnya.

Penyiapannya dilakukan sekali, seperti pada Listing \ref{lst:c6-siapkan}.

\begin{lstlisting}[language=bash, caption={Menyiapkan basis data, data uji, dan dependensi, memakai \berkas{sources/chapter06/siapkan-data.sh}}, label={lst:c6-siapkan}]
cd sources/chapter06
docker compose up -d
./siapkan-data.sh
python3 -m venv ../.venv
../.venv/bin/pip install "psycopg[binary]" pytest hypothesis jsonschema \
  coverage ruff
npm install
\end{lstlisting}

\ifajar
  \subsection*{Latihan 1: Membandingkan Tiga Lapisan Pengujian}
\else
  \subsection*{Prosedur 1: Membandingkan Tiga Lapisan Pengujian}
\fi

Latihan ini membuktikan bahwa ketiga lapisan berbeda harganya, dan menangkap
cacat yang berbeda pula.

\begin{enumerate}
  \item Jalankan ketiga lapisan seperti pada Listing
        \ref{lst:c6-lat1-jalan}, lalu catat jumlah pengujian dan waktunya ke
        Tabel \ref{tab:c6-lat1-hasil}. Tabel \ref{tab:c6-lat1-contoh}
        memperlihatkan contoh yang sudah terisi.
  \item Hitung berapa kali lipat satu pengujian E2E dibandingkan rata-rata
        satu unit test.
  \item Matikan aplikasinya, lalu jalankan ulang lapisan E2E dan catat apa
        yang terjadi. Nyalakan kembali sesudahnya.
  \item Jawab tiga pertanyaan berikut. Pertama, cacat apa yang hanya
        tertangkap lapisan E2E? Kedua, mengapa unit test tetap dibuat paling
        banyak walaupun paling sempit jangkauannya? Ketiga, apa akibatnya
        bila seluruh pengujian dipindahkan ke lapisan E2E?
\end{enumerate}

\begin{lstlisting}[language=bash, caption={Menjalankan ketiga lapisan pengujian, memakai \berkas{sources/chapter06/}}, label={lst:c6-lat1-jalan}]
source ../.venv/bin/activate
pytest -q test_harga.py test_properti.py test_double.py
pytest -q test_integrasi.py test_kontrak.py
python aplikasi.py &
npx playwright test --config e2e/playwright.config.ts
\end{lstlisting}

\begin{table}[!htb]
  \centering
  \caption{Contoh hasil Latihan 1 yang sudah terisi}
  \label{tab:c6-lat1-contoh}
  \begin{tabularx}{\textwidth}{lrrX}
    \hline
    \textbf{Lapisan} & \textbf{Jumlah} & \textbf{Waktu} & \textbf{Catatan} \\
    \hline
    Unit & 19 & 0,99 s & Tanpa basis data dan tanpa jaringan \\
    \hline
    Integrasi dan kontrak & 15 & 1,27 s & Basis data di container \\
    \hline
    E2E & 5 & 2,90 s & Chrome sungguhan, satu \textit{worker} \\
    \hline
  \end{tabularx}
\end{table}

\begin{table}[!htb]
  \centering
  \caption{Lembar isian hasil pengukuran Latihan 1}
  \label{tab:c6-lat1-hasil}
  \begin{tabularx}{\textwidth}{lrrX}
    \hline
    \textbf{Lapisan} & \textbf{Jumlah} & \textbf{Waktu} & \textbf{Catatan} \\
    \hline
    Unit & \dots & \dots & \dots \\
    \hline
    Integrasi dan kontrak & \dots & \dots & \dots \\
    \hline
    E2E & \dots & \dots & \dots \\
    \hline
    E2E tanpa aplikasi menyala & \dots & \dots & \dots \\
    \hline
  \end{tabularx}
\end{table}

\ifajar
  \subsection*{Latihan 2: Coverage yang Menipu}
\else
  \subsection*{Prosedur 2: Coverage yang Menipu}
\fi

Latihan ini membuktikan bahwa \textit{coverage} 100 persen tidak menjamin
cacat tertangkap.

\begin{enumerate}
  \item Ukur \textit{coverage} \berkas{harga.py} untuk seluruh
        \berkas{test\_harga.py}, lalu ukur lagi tanpa pengujian batas,
        seperti pada Listing \ref{lst:c6-lat2-jalan}.
  \item Ubah \texttt{>=} menjadi \texttt{>} pada \texttt{hitung\_diskon} di
        \berkas{harga.py}. Ini disebut mutasi.
  \item Jalankan kedua suite terhadap kode yang sudah dimutasi, lalu catat
        mana yang lolos dan mana yang gagal ke Tabel
        \ref{tab:c6-lat2-hasil}. Tabel \ref{tab:c6-lat2-contoh}
        memperlihatkan contoh yang sudah terisi.
  \item Kembalikan \berkas{harga.py} seperti semula, lalu pastikan seluruh
        pengujian kembali hijau.
  \item Jawab dua pertanyaan berikut. Pertama, mengapa \textit{coverage}
        tetap 100 persen padahal satu pengujian dibuang? Kedua, pengujian
        seperti apa yang perlu ditambahkan agar mutasi semacam ini selalu
        tertangkap?
\end{enumerate}

\begin{lstlisting}[language=bash, caption={Mengukur coverage dua kali dan menjalankan mutasi, memakai \berkas{sources/chapter06/harga.py}}, label={lst:c6-lat2-jalan}]
source ../.venv/bin/activate
coverage run -m pytest -q test_harga.py
coverage report --include="harga.py"
coverage run -m pytest -q test_harga.py -k "not batas_grosir"
coverage report --include="harga.py"

# Keluarannya:
# Name       Stmts   Miss  Cover
# harga.py      23      0   100%
\end{lstlisting}

\begin{table}[!htb]
  \centering
  \caption{Contoh hasil Latihan 2 yang sudah terisi}
  \label{tab:c6-lat2-contoh}
  \begin{tabularx}{\textwidth}{Xllr}
    \hline
    \textbf{Suite} & \textbf{Coverage} & \textbf{Sesudah mutasi} & \textbf{Gagal} \\
    \hline
    Seluruh \berkas{test\_harga.py} & 100\% & tertangkap & 1 \\
    \hline
    Tanpa pengujian batas & 100\% & lolos & 0 \\
    \hline
  \end{tabularx}
\end{table}

\begin{table}[!htb]
  \centering
  \caption{Lembar isian hasil pengukuran Latihan 2}
  \label{tab:c6-lat2-hasil}
  \begin{tabularx}{\textwidth}{Xllr}
    \hline
    \textbf{Suite} & \textbf{Coverage} & \textbf{Sesudah mutasi} & \textbf{Gagal} \\
    \hline
    Seluruh \berkas{test\_harga.py} & \dots & \dots & \dots \\
    \hline
    Tanpa pengujian batas & \dots & \dots & \dots \\
    \hline
    Mutasi pilihan sendiri & \dots & \dots & \dots \\
    \hline
  \end{tabularx}
\end{table}

\ifajar
  \subsection*{Latihan 3: Menangkap N+1 Lewat Jumlah Query}
\else
  \subsection*{Prosedur 3: Menangkap N+1 Lewat Jumlah Query}
\fi

Latihan ini membuktikan bahwa dua alamat yang memberi \textit{response} sama
persis dapat berbeda jauh biayanya, dan bahwa bedanya hanya terbaca bila
jumlah \textit{query}-nya ikut diukur.

\begin{enumerate}
  \item Nyalakan aplikasinya, lalu panggil kedua alamat dengan
        \texttt{curl} seperti pada Listing \ref{lst:c6-lat3-jalan}. Catat
        header \texttt{X-Query-Count} keduanya.
  \item Hitung berapa item yang dimiliki pesanan tersebut, lalu cocokkan
        dengan rumus dua ditambah jumlah item.
  \item Bandingkan isi kedua \textit{response} dengan \texttt{diff}. Catat
        apakah ada bedanya.
  \item Data uji Bab 3 memberi tiap pesanan tepat tiga item, sehingga jumlah
        item yang berbeda dibuat sendiri. Tambahkan dua item ke pesanan 2,
        ulangi ketiga langkah di atas untuk pesanan itu, lalu hapus kembali
        kedua item tambahannya, seperti pada Listing
        \ref{lst:c6-lat3-tambah}.
  \item Isi Tabel \ref{tab:c6-lat3-hasil}. Tabel \ref{tab:c6-lat3-contoh}
        memperlihatkan contoh yang sudah terisi.
  \item Jawab dua pertanyaan berikut. Pertama, mengapa pengujian yang hanya
        membandingkan isi \textit{response} tidak menangkap cacat ini? Kedua,
        mengapa jumlah \textit{query} lebih layak dijadikan batas di
        \textit{Continuous Integration} daripada waktu jalannya?
\end{enumerate}

\begin{lstlisting}[language=bash, caption={Membaca jumlah query dari header, tersedia di \berkas{sources/chapter06/aplikasi.py}}, label={lst:c6-lat3-jalan}]
source ../.venv/bin/activate
python aplikasi.py &
curl -s -D - -o sehat.json http://127.0.0.1:8020/api/pesanan/1 \
  | grep -i x-query-count
curl -s -D - -o n1.json http://127.0.0.1:8020/api/pesanan-n1/1 \
  | grep -i x-query-count
diff sehat.json n1.json && echo "isi response sama persis"

# Keluarannya:
# X-Query-Count: 2
# X-Query-Count: 5
# isi response sama persis
\end{lstlisting}

\begin{lstlisting}[language=bash, caption={Menambah dua item ke pesanan 2 lalu mengukurnya, memakai \berkas{sources/chapter06/aplikasi.py}}, label={lst:c6-lat3-tambah}]
docker compose exec -T db psql -U admin -d toko -c \
  "INSERT INTO item_pesanan (pesanan_id, buku_id, jumlah, harga_satuan)
   SELECT 2, id, 1, harga FROM buku WHERE id IN (101, 102)"
curl -s -D - -o sehat2.json http://127.0.0.1:8020/api/pesanan/2 \
  | grep -i x-query-count
curl -s -D - -o n1-2.json http://127.0.0.1:8020/api/pesanan-n1/2 \
  | grep -i x-query-count
diff sehat2.json n1-2.json && echo "isi response sama persis"
docker compose exec -T db psql -U admin -d toko -c \
  "DELETE FROM item_pesanan WHERE pesanan_id = 2 AND buku_id IN (101, 102)"

# Keluarannya:
# INSERT 0 2
# X-Query-Count: 2
# X-Query-Count: 7
# isi response sama persis
# DELETE 2
\end{lstlisting}

\begin{table}[!htb]
  \centering
  \caption{Contoh hasil Latihan 3 yang sudah terisi}
  \label{tab:c6-lat3-contoh}
  \begin{tabularx}{\textwidth}{lrrrX}
    \hline
    \textbf{Pesanan} & \textbf{Item} & \textbf{Query sehat} & \textbf{Query N+1} & \textbf{Isi response} \\
    \hline
    1 & 3 & 2 & 5 & sama persis \\
    \hline
    2, ditambah dua item & 5 & 2 & 7 & sama persis \\
    \hline
  \end{tabularx}
\end{table}

\begin{table}[!htb]
  \centering
  \caption{Lembar isian hasil pengukuran Latihan 3}
  \label{tab:c6-lat3-hasil}
  \begin{tabularx}{\textwidth}{lrrrX}
    \hline
    \textbf{Pesanan} & \textbf{Item} & \textbf{Query sehat} & \textbf{Query N+1} & \textbf{Isi response} \\
    \hline
    1 & \dots & \dots & \dots & \dots \\
    \hline
    2, ditambah dua item & \dots & \dots & \dots & \dots \\
    \hline
  \end{tabularx}
\end{table}

\ifajar
  \subsection*{Latihan 4: Menegakkan Quality Gate}
\else
  \subsection*{Prosedur 4: Menegakkan Quality Gate}
\fi

Latihan ini memasang seluruh pengukuran di atas sebagai satu \textit{gate} yang
\textit{verdict}-nya dapat dibaca mesin.

\begin{enumerate}
  \item Baca \berkas{sources/chapter06/mutu.json}, lalu tulis dugaan
        pemeriksaan mana yang akan dilanggar dan mengapa.
  \item Jalankan \textit{gate}-nya seperti pada Listing \ref{lst:c6-lat4-jalan},
        lalu catat \textit{verdict} tiap baris beserta status keluarnya ke
        Tabel \ref{tab:c6-lat4-hasil}. Tabel \ref{tab:c6-lat4-contoh}
        memperlihatkan contoh yang sudah terisi.
  \item Longgarkan batas \textit{query} alamat N+1 menjadi tiga puluh, lalu
        jalankan ulang dan catat status keluarnya.
  \item Kembalikan batasnya menjadi dua, lalu perbaiki cacatnya dengan
        mengarahkan alamat \texttt{/api/pesanan-n1/} ke jalur sehat, yaitu
        mengganti \texttt{n\_plus\_1=True} menjadi \texttt{n\_plus\_1=False}
        di \berkas{sources/chapter06/aplikasi.py}. Jalankan ulang
        \textit{gate}-nya, lalu catat pemeriksaan mana yang kini dilanggar.
  \item Sesuaikan pengujian di \berkas{sources/chapter06/test\_kontrak.py}
        yang masih mengharapkan jumlah \textit{query} lama: harapannya diubah
        menjadi \texttt{QUERY\_JALUR\_SEHAT}, dan baris \texttt{jumlah\_item}
        yang tidak terpakai lagi dihapus. Jalankan ulang \textit{gate}-nya,
        catat status keluarnya, lalu kembalikan kedua file itu seperti semula
        sebelum lanjut ke Latihan 5.
  \item Jawab dua pertanyaan berikut. Pertama, apa bedanya melonggarkan batas
        dengan memperbaiki cacatnya, dilihat dari status keluarnya? Kedua,
        mengapa \textit{Continuous Integration} membaca status keluar dan
        bukan tulisan di layarnya?
\end{enumerate}

\begin{lstlisting}[language=bash, caption={Menjalankan gate mutu dan membaca status keluarnya, tersedia di \berkas{sources/chapter06/gate-mutu.py}}, label={lst:c6-lat4-jalan}]
source ../.venv/bin/activate
python gate-mutu.py
echo $?

# Keluarannya:
# Gagal: 1 batas mutu dilanggar.
# 1
\end{lstlisting}

\begin{table}[!htb]
  \centering
  \caption{Contoh hasil Latihan 4 yang sudah terisi}
  \label{tab:c6-lat4-contoh}
  \begin{tabularx}{\textwidth}{Xllr}
    \hline
    \textbf{Pemeriksaan} & \textbf{Batas} & \textbf{Terukur} & \textbf{Verdict} \\
    \hline
    waktu unit & 10 s & 1,89 s & lolos \\
    \hline
    coverage \berkas{harga.py} & 95\% & 100\% & lolos \\
    \hline
    temuan \texttt{ruff} & 0 & 0 & lolos \\
    \hline
    query \texttt{/api/pesanan-n1/1} & 2 & 5 & LANGGAR \\
    \hline
    Status keluar & 0 & 1 & LANGGAR \\
    \hline
    \texttt{pytest} integrasi, sesudah alamat diarahkan & 0 & 1 & LANGGAR \\
    \hline
    Status keluar sesudah diperbaiki & 0 & 0 & lolos \\
    \hline
  \end{tabularx}
\end{table}

Baris keenam pada contoh terisi memperlihatkan akibat yang sering terlupa.
Begitu alamatnya diarahkan ke jalur sehat, satu pengujian kontrak ikut gagal,
karena pengujian itu masih mengunci jumlah \textit{query} jalur N+1 yang lama.
Pengujian semacam itu diperbarui bersama perbaikannya, dan baru sesudahnya
status keluarnya menjadi 0.

\begin{table}[!htb]
  \centering
  \caption{Lembar isian hasil pengukuran Latihan 4}
  \label{tab:c6-lat4-hasil}
  \begin{tabularx}{\textwidth}{Xllr}
    \hline
    \textbf{Pemeriksaan} & \textbf{Batas} & \textbf{Terukur} & \textbf{Verdict} \\
    \hline
    waktu unit & \dots & \dots & \dots \\
    \hline
    waktu integrasi & \dots & \dots & \dots \\
    \hline
    coverage \berkas{harga.py} & \dots & \dots & \dots \\
    \hline
    temuan \texttt{ruff} & \dots & \dots & \dots \\
    \hline
    query \texttt{/api/pesanan/1} & \dots & \dots & \dots \\
    \hline
    query \texttt{/api/pesanan-n1/1} & \dots & \dots & \dots \\
    \hline
    Status keluar sebelum dan sesudah dilonggarkan & \dots & \dots & \dots \\
    \hline
    Pemeriksaan yang dilanggar sesudah alamat diarahkan & \dots & \dots & \dots \\
    \hline
    Status keluar sesudah cacatnya diperbaiki & \dots & \dots & \dots \\
    \hline
  \end{tabularx}
\end{table}

\ifajar
  \subsection*{Latihan 5: Membedah Pengujian E2E}
\else
  \subsection*{Prosedur 5: Membedah Pengujian E2E}
\fi

Latihan ini membuktikan bahwa lapisan E2E paling bergoyang waktunya, dan
memperlihatkan jejak yang ditinggalkannya ketika sebuah pengujian gagal.

\begin{enumerate}
  \item Jalankan lapisan E2E tiga kali berurutan seperti pada Listing
        \ref{lst:c6-lat5-jalan}, lalu catat waktu tiap kali jalan ke Tabel
        \ref{tab:c6-lat5-hasil}. Tabel \ref{tab:c6-lat5-contoh}
        memperlihatkan contoh yang sudah terisi.
  \item Jalankan sekali lagi dengan \texttt{-{}-repeat-each=3}, lalu catat
        jumlah pengujian dan waktunya.
  \item Ubah \texttt{QUERY\_JALUR\_SEHAT} di \berkas{sources/chapter06/e2e/pesanan.spec.ts}
        menjadi 3, jalankan ulang dengan \texttt{-{}-trace on}, lalu catat
        pesan gagalnya dan isi direktori \berkas{test-results/}.
  \item Kembalikan angkanya menjadi 2, lalu pastikan kelima pengujian kembali
        hijau.
  \item Jawab dua pertanyaan berikut. Pertama, mengapa waktu lapisan E2E
        bergoyang jauh lebih lebar daripada lapisan unit? Kedua, mengapa
        \textit{gate} mutu menilai waktu lapisan unit dan integrasi, tetapi
        tidak menilai waktu lapisan E2E?
\end{enumerate}

\begin{lstlisting}[language=bash, caption={Menjalankan lapisan E2E berulang kali beserta jejaknya, memakai \berkas{sources/chapter06/e2e/playwright.config.ts}}, label={lst:c6-lat5-jalan}]
npx playwright test --config e2e/playwright.config.ts
npx playwright test --config e2e/playwright.config.ts --repeat-each=3
npx playwright test --config e2e/playwright.config.ts --trace on
du -sh test-results

# Keluarannya:
# 5 passed (2.9s)
# 15 passed (6.7s)
# 5 passed (2.8s)
# 369K  test-results
\end{lstlisting}

\begin{table}[!htb]
  \centering
  \caption{Contoh hasil Latihan 5 yang sudah terisi}
  \label{tab:c6-lat5-contoh}
  \begin{tabularx}{\textwidth}{Xrr}
    \hline
    \textbf{Cara menjalankan} & \textbf{Pengujian} & \textbf{Waktu} \\
    \hline
    Biasa, jalan pertama & 5 & 2,1 s \\
    \hline
    Biasa, jalan kedua & 5 & 2,9 s \\
    \hline
    Biasa, jalan ketiga & 5 & 5,6 s \\
    \hline
    Dengan \texttt{-{}-repeat-each=3} & 15 & 6,7 s \\
    \hline
  \end{tabularx}
\end{table}

Ketiga angka pada tiga baris pertama diambil apa adanya, termasuk 5,6 detik
yang jauh lebih besar daripada dua lainnya. Sebabnya browser sungguhan ikut
berebut prosesor dengan proses lain di komputer yang sama, dan itulah alasan
waktu E2E tidak layak dipakai sebagai batas di \textit{Continuous
Integration}.

\begin{table}[!htb]
  \centering
  \caption{Lembar isian hasil pengukuran Latihan 5}
  \label{tab:c6-lat5-hasil}
  \begin{tabularx}{\textwidth}{Xrr}
    \hline
    \textbf{Cara menjalankan} & \textbf{Pengujian} & \textbf{Waktu} \\
    \hline
    Biasa, jalan pertama & \dots & \dots \\
    \hline
    Biasa, jalan kedua & \dots & \dots \\
    \hline
    Biasa, jalan ketiga & \dots & \dots \\
    \hline
    Dengan \texttt{-{}-repeat-each=3} & \dots & \dots \\
    \hline
    Sesudah angka harapan diubah menjadi 3 & \dots & \dots \\
    \hline
  \end{tabularx}
\end{table}

Setelah seluruh latihan selesai, container dihentikan dengan
\texttt{docker compose down -v}. Opsi \texttt{-v} ikut menghapus volume
datanya, sehingga penyiapan berikutnya dimulai dari keadaan bersih.
